\section{Method-Based Policies}\label{sec:policies}

By Proposition~\ref{prop:fm-vs-mp}(ii) the weak column for method-based policies coincides with the
weak column for fixed methods, and by Proposition~\ref{prop:fm-vs-mp}(i) the primitive rows
coincide with strong OD (Theorem~\ref{thm:prim-po}). We therefore concentrate on strong and
(fair) strong cyclic policies for problems with a hierarchy.

\subsection{Upper bounds}

All upper bounds come from the games of Section~\ref{sec:games}, played with progression only
(Lemma~\ref{lem:delay}). What varies between the classes is the size of the networks that
progression can reach \parencite{alford:et:al:icaps-15}.

\begin{theorem}[\CBb, Thms.~1--5]\label{thm:mp-upper}
  Strong and fair strong cyclic method-based plan existence is
  \begin{itemize}
    \item in \EXPTIME{} for acyclic totally ordered, regular (\REG{} and \REGtn), and
      tail-recursive totally ordered problems;
    \item in \EXPSPACE{} for acyclic partially ordered problems;
    \item in \TWOEXP{} for tail-recursive partially ordered problems.
  \end{itemize}
\end{theorem}

\begin{proof}
  For acyclic problems every play of the game of Theorem~\ref{thm:mp-andor} is finite, and strong
  and fair solutions coincide (Proposition~\ref{prop:fair-basic}). Under total order the networks
  reachable by progression have polynomial size, so the game is decided in alternating polynomial
  space, i.e.\ \EXPTIME. Under partial order they have exponential size and plays have exponential
  length, which is alternating exponential time, i.e.\ \EXPSPACE.

  For the other classes plays can be infinite, and we compute the winning region explicitly. In a
  regular problem a reachable network consists of a subset of the primitive tasks of one method
  (or of $\tnI$) plus possibly the compound task, so there are exponentially many networks up to
  isomorphism. For tail-recursive problems progression keeps networks polynomially large under total
  order and exponentially large under partial order. Together with the $2^{|F|}$ states, the game
  graph is exponential, respectively doubly exponential. The attractor for strong solutions and the
  nested fixpoint of Theorem~\ref{thm:mp-fair-game} for fair ones are computed in time polynomial in
  the size of the graph.
\end{proof}

We note that \CBb{} also suggest deciding strong cyclic solutions with their alternating
progression procedure, after dropping its cycle check. This computes strong rather than strong
cyclic solutions, since an alternating machine rejects infinite branches. The upper bounds for
strong cyclic solutions therefore rest on the explicit graph computation, as above.

\subsection{\EXPTIME{} lower bounds}\label{sec:exptime-hard}

\CBb{} prove \EXPTIME-hardness for acyclic totally ordered problems (Thm.~1) by simulating an
alternating polynomial-space machine with an exponentially long sequence of step tasks. The
construction has no acceptance check: the method that lets the simulation stop can be used at any
time, so every instance has a trivial solution. Their lower bound for regular problems (Thm.~3)
encodes FOND STRIPS planning and needs a method that the proof does not mention. We give one
construction that covers all three \EXPTIME-complete classes.

\begin{theorem}\label{thm:exptime-hard}
  Strong and fair strong cyclic method-based plan existence is \EXPTIME-hard for acyclic totally
  ordered problems and for regular totally ordered problems (hence also for regular partially
  ordered and tail-recursive totally ordered problems).
\end{theorem}

The source problem is deterministic plan existence for totally ordered HTN problems, which is
\EXPTIME-complete \parencite{alford:et:al:icaps-15}. Such a problem is decided by a dynamic
program over triples $(v,x,v')$: ``task $x$ can be refined into an executable sequence leading
from state $v$ to state $v'$''. A primitive $x$ needs a single applicable step, and a compound $x$
needs a method $x\to x_1\cdots x_k$ and intermediate states $v=u_0,u_1,\dots,u_k=v'$ such that all
triples $(u_{i-1},x_i,u_i)$ hold. The table has exponentially many triples, and it is the least
fixpoint of this rule.

\begin{proof}[Proof sketch]
  We build a FOND problem whose state stores a \emph{current triple} $(v,x,v')$ in polynomially many
  Boolean fields. The policy proves that the triple is in the table by walking down a derivation
  tree, with the environment choosing the branch. One \emph{round} of this walk is a small program.
  \begin{itemize}
    \item For a compound $x$, the policy chooses a method $x\to x_1\cdots x_k$ of the source and
      guesses the intermediate states $u_1,\dots,u_{k-1}$, bit by bit, by choosing among
      alternative methods. The program checks $u_0=v$ and $u_k=v'$. Then a nondeterministic action
      with $k$ outcomes chooses a child $i$, and the current triple becomes $(u_{i-1},x_i,u_i)$.
    \item For a primitive $x$, the program checks bit by bit that the action is applicable in $v$
      and leads to $v'$, and it sets a flag $\mathit{done}$.
  \end{itemize}
  Programs are compiled into the hierarchy: the compound task ``program $r$ at instruction $pc$''
  has one method per alternative of the instruction, namely the corresponding action followed by
  ``program $r$ at instruction $pc+1$''. This is totally ordered and regular.

  Rounds are supplied in one of two ways. In the \emph{acyclic} variant a doubling clock
  $\mathit{clk}_{j+1}\to\mathit{clk}_j\,\mathit{clk}_j$ produces $2^n$ rounds, enough for any
  derivation path, and a final task requires $\mathit{done}$. A round with $\mathit{done}$ set may
  do nothing. In the \emph{regular} variant a compound task $\mathit{round}$ has a method for every
  program, each ending in $\mathit{round}$ again, and a method $\mathit{round}\to\mathit{fin}$.

  If the initial triple (initial state, initial task, any state) is in the table, the policy plays
  along derivations and wins in every branch. Conversely, an invariant shows that every reachable
  configuration is either done or has its triple in the table: a round can only be won against all
  children if all children's triples hold. For strong policies this is all. In the regular variant
  a fair policy could keep looping through rounds whose guesses are consistent but never reach the
  bottom of a derivation. To exclude this we add a round counter that bounds the number of rounds by
  the depth of the dynamic program. Then every play is finite, and fair solutions are strong.
\end{proof}

The constructed problems use negative preconditions. Actions have as many outcomes as the source
has subtasks in a method.

\subsection{Tail-recursive partially ordered problems}\label{sec:tr-po}

\begin{theorem}\label{thm:tr-po}
  Strong and fair strong cyclic method-based plan existence for tail-recursive partially ordered
  problems is \TWOEXP-complete.
\end{theorem}

Membership is Theorem~\ref{thm:mp-upper}. For strong policies \CBb{} (Thm.~5) claim hardness with
the remark that the \EXPSPACE-hardness proof for deterministic problems can be extended to
alternating machines ``in the same fashion'' as their Thm.~2; for strong cyclic policies they leave
hardness open. We give a complete reduction that works for both.

\paragraph{The source.} Every language in \TWOEXP{} is decided by a deterministic machine running in
time $2^{2^{e(n)}}$, and acceptance can be phrased as a \emph{tableau game}. The existential player
claims ``at time $t$, cell $i$ holds symbol $a$''. To support a claim with $t>0$ she names the
contents of the three cells $i-1,i,i+1$ at time $t-1$ that determine it, and the universal player
picks one of them to continue with. A claim at $t=0$ is checked against the input. The numbers
$t$ and $i$ have $B=2^{e(n)}$ bits, exponentially many, so a claim does not fit into a state.

\paragraph{Claims on a tape of threads.}
We store a claim in the task network, as two interleaved $B$-bit counters on a ``tape'' of
$4B$ cells. One round of the game sweeps this tape once. A round is played by a loop
\[
  \mathit{loop}\to\mathit{block}\cdot\mathit{loop},
\]
which is tail-recursive because the recursion is in the last task. Two such loops alternate as
writer and reader. A block expands, through acyclic doubling with unordered methods, into $4B$
\emph{threads}, one per cell. A thread writes the address of its cell, reads the bit that the
corresponding thread of the previous round left behind, applies one step of a \emph{transducer} to
it, and leaves the new bit for the next round. The transducer decrements $t$ (with a borrow that
travels with the sweep) and moves $i$ by the offset $-1$, $0$ or $+1$ that the environment chose
at the end of the previous round. The bit left behind is the tail of the thread, which is
unordered with respect to the next block, so the next round can read it whenever its own thread for
that cell runs. At the end of a round the environment chooses the next offset, and the policy names
the claimed symbol. When $t$ reaches $0$, the first cell guesses which interval of the first row
(input, padding, blanks) contains $i$, and the sweep verifies the guess by comparing $i$ with the
interval bounds, which have polynomially many set bits.

\begin{proof}[Proof sketch]
  The construction is computable in polynomial time: the interval bounds are exponentially long
  numbers, but they are accessed only bitwise at polynomially many positions.

  Soundness: every round of a strong (or fair) policy realises a correct move of the tableau game,
  and the time decreases in every round. By induction on the time, a won claim is a true entry
  of the tableau. Completeness: the canonical strategy of the tableau game, executed by writing the
  claimed symbols and following the environment's offsets, wins. For fair policies we only use
  that a fair policy cannot leave the set of positions it has committed to (a safety property). The
  decrease of the time takes care of termination, so the same argument applies.
\end{proof}

\subsection{Totally ordered problems with an arbitrary hierarchy}\label{sec:to-arb}

Totally ordered problems with arbitrary recursion are missing from the table of \CBb{}. For
deterministic problems they are \EXPTIME-complete \parencite{alford:et:al:icaps-15}. With
policies they are one exponential harder.

\begin{theorem}\label{thm:to-arb}
  Strong and fair strong cyclic method-based plan existence for totally ordered problems is
  \TWOEXP-complete.
\end{theorem}

\paragraph{Upper bound for strong policies.}
The decisive tool is sequential composition. Write $u;v$ for a network in which every task of $u$
precedes every task of $v$, and say that $u$ is won from $s$ \emph{with target} $S'$ if the planner
can execute $u$ completely against all outcomes so that every final state lies in $S'$.

\begin{lemma}[composition]\label{lem:comp}
  $u;v$ is won from $s$ iff, for some set $S'$, $u$ is won from $s$ with target $S'$ and $v$ is won
  from every state in $S'$.
\end{lemma}

\begin{proof}
  Since no task of $v$ can be touched before $u$ is finished, a strategy for $u;v$ first plays $u$
  to completion and then continues. Take $S'$ to be the set of states in which it finishes $u$.
  Conversely, play a strategy for $u$ and then, from the reached state, one for $v$.
\end{proof}

For a totally ordered network $c\,w$ that starts with a compound task, the lemma reduces the game to
\emph{summaries} $(c,s,S')$: ``$c$, started in $s$, can be completed with target $S'$''. A summary
holds iff some method $c\to x_1\cdots x_k$ can be executed from $s$ with target $S'$, and executing a
word of tasks again splits at its first compound task into a summary and a rest. This yields an
AND/OR game whose configurations, a task, a state, a remainder of a method network and an explicit
set $S'$ of states, have exponential size. Its attractor is computed in doubly exponential time.
The game is the least fixpoint, which matches strong solutions: a winning strategy has finite height,
so the summaries used in it are well founded. Turning a derivation of summaries back into a policy
needs a rank argument, so that the policy never returns to an isomorphic node. This is the
``rank'' that the unpublished notes in which the table was proposed leave open.

\paragraph{Upper bound for fair policies.}
Composition holds for fair policies as well, but the obvious summary table, a greatest instead of a
least fixpoint, is wrong. A fair policy for $c\,w$ must satisfy two requirements of different
kinds. Locally, every node inside the execution of $c$ must be able to reach \emph{some} exit of
$c$. And for the enclosing network to make progress, the start of $c$ must be able to reach a
\emph{particular} exit from which the rest can be completed. Neither the least nor the greatest
fixpoint of the summary table captures both. We therefore use a summary game with two progress
modes: a position records, besides the remaining word, the state and the allowed exit set, whether
it must reach \emph{some} exit or a \emph{specific} exit $e$. A call of a compound task splits into
the call itself (with its own exit set and progress target) and the continuation. Safety is
required for all exits of the call, progress only through the chosen target. The game is a fair
game with separate successor sets for safety and progress, its configurations have exponential
size, and its nested fixpoint is computed in doubly exponential time.

\paragraph{Lower bound.}
The unpublished notes contain a sketch of the hardness proof, which leaves two families of
compound tasks undefined, misses the base case of the induction and does not show that the planner
cannot profit from inconsistent encodings. We give a complete construction along the same lines.

\begin{proof}[Proof sketch of hardness]
  As in Section~\ref{sec:tr-po} we use the tableau of a deterministic machine that runs in time
  $2^D$ with $D=2^b$, $b$ polynomial, on a doubly exponentially padded input. A claim
  $\Claim(t,i,a)$ involves two numbers with $D$ binary \emph{digits}. A digit index has only $b$
  bits and fits into the state, but the numbers do not, and in a totally ordered network there is
  no room for parallel threads either.

  The key idea is \emph{digit challenges}. The claim task for time $t$ is refined, left-recursively,
  into the claim for time $t-1$ followed by a \emph{verification chain}. The chain runs through all
  $D$ digit positions, and for each position it executes a small automaton step that checks one
  local relation between the digits of the old and new numbers, e.g.\ ``$t$ is one more than the
  inner time'', with the carry handled by the automaton. The types of these steps, and hence the
  numbers themselves, are fixed when the policy chooses the methods of the chain. Recursion
  bottoms out at time $0$, which is executed \emph{first} because it is leftmost. There the
  environment chooses a digit index $\mathit{Check}$, once and for all. From then on the state
  records only digit $\mathit{Check}$ of each number, and every verification step compares the
  digit it is about against $\mathit{Check}$ and is vacuous for the other digits.

  Because $\mathit{Check}$ is chosen by the environment, after the verification chains have been
  committed, a strong policy must be consistent at \emph{every} digit, so the committed numbers are
  well defined and satisfy the relations. Which of the three predecessor cells is verified is again
  chosen by the environment, as a \emph{mode}, before the inner claim is refined. The policy may then
  choose the inner claim's methods depending on the mode, which is exactly what method-based
  policies allow. By induction on $t$, $\Claim(t,i,a)$ can be won iff the tableau has $a$ at
  $(t,i)$.

  The same construction works for fair policies, but the induction needs care, because a fair
  policy may recurse forever along some challenges. One shows first, by induction on the rank of
  the attractor, that every claim the policy commits to covers all challenged digits. Coverage forces
  the time to decrease, so the claims are sound by induction on the time.
\end{proof}
